%-------------------------
% Resume in Latex
% Author : Gyanig Kumar
% Based off of: https://github.com/sb2nov/resume
% License : MIT
%------------------------

\documentclass[letterpaper,10pt]{article}

\usepackage{latexsym}
\usepackage[empty]{fullpage}
\usepackage{titlesec}
\usepackage{marvosym}
\usepackage[usenames,dvipsnames]{color}
\usepackage{verbatim}
\usepackage{enumitem}
\usepackage{hyperref}
\usepackage{fancyhdr}
\usepackage[english]{babel}
\usepackage{tabularx}
\usepackage{fontawesome5}
\usepackage{multicol}
\usepackage{academicons}
\setlength{\multicolsep}{-3.0pt}
\setlength{\columnsep}{-1pt}
\input{glyphtounicode}
\usepackage{xcolor}
\usepackage[none]{hyphenat}
\hypersetup{
  colorlinks = true,% hyperlinks will be black
  filecolor=blue,
  linkcolor = blue,% hyperlink borders will be red
  urlcolor=black,
  pdfborderstyle={/S/U/W 1}% border style will be underline of width 1pt
}


%----------FONT OPTIONS----------
% sans-serif
% \usepackage[sfdefault]{FiraSans}
% \usepackage[sfdefault]{roboto}
% \usepackage[sfdefault]{noto-sans}
% \usepackage[default]{sourcesanspro}

% serif
% \usepackage{CormorantGaramond}
% \usepackage{charter}


\pagestyle{fancy}
\fancyhf{} % clear all header and footer fields
\fancyfoot{}
\renewcommand{\headrulewidth}{0pt}
\renewcommand{\footrulewidth}{0pt}

% Adjust margins
\addtolength{\oddsidemargin}{-0.5in}
% \addtolength{\evensidemargin}{-0.2in}
\addtolength{\textwidth}{1.0in}
\addtolength{\topmargin}{-.8in}
\addtolength{\textheight}{1.4in}

\urlstyle{same}

\raggedbottom
\raggedright
\setlength{\tabcolsep}{0in}

% Sections formatting
\titleformat{\section}{
  \vspace{-4pt}\scshape\raggedright\large\bfseries
}{}{0em}{}[\color{black}\titlerule \vspace{-5pt}]

% Ensure that generate pdf is machine readable/ATS parsable
\pdfgentounicode=1

%-------------------------
% Custom commands
\newcommand{\resumeItem}[1]{
  \item\small{
    {#1 \vspace{-2pt}}
  }
}

\newcommand{\classesList}[4]{
    \item\small{
        {#1 #2 #3 #4 \vspace{-2pt}}
  }
}

\newcommand{\resumeSubheading}[4]{
  \vspace{-2pt}\item
    \begin{tabular*}{1.0\textwidth}[t]{l@{\extracolsep{\fill}}r}
      \textbf{#1} & \textbf{\small #2} \\
      \textit{\small#3} & \textit{\small #4} \\
    \end{tabular*}\vspace{-7pt}
}

\newcommand{\resumeSubSubheading}[2]{
    \item
    \begin{tabular*}{0.97\textwidth}{l@{\extracolsep{\fill}}r}
      \textit{\small#1} & \textit{\small #2} \\
    \end{tabular*}\vspace{-7pt}
}

\newcommand{\resumeProjectHeading}[2]{
    \item
    \begin{tabular*}{1.001\textwidth}{l@{\extracolsep{\fill}}r}
      \small#1 & \textbf{\small #2}\\
    \end{tabular*}\vspace{-7pt}
}

\newcommand{\resumeSubItem}[1]{\resumeItem{#1}\vspace{-4pt}}

\renewcommand\labelitemi{$\vcenter{\hbox{\tiny$\bullet$}}$}
\renewcommand\labelitemii{$\vcenter{\hbox{\tiny$\bullet$}}$}

\newcommand{\resumeSubHeadingListStart}{\begin{itemize}[leftmargin=0.0in, label={}]}
\newcommand{\resumeSubHeadingListEnd}{\end{itemize}}
\newcommand{\resumeItemListStart}{\begin{itemize}}
\newcommand{\resumeItemListEnd}{\end{itemize}\vspace{-5pt}}
% \newcommand*{\scholarsocialsymbol}{\includegraphics[height=.7\baselineskip]{google-scholar.pdf}}
%-------------------------------------------
%%%%%%  RESUME STARTS HERE  %%%%%%%%%%%%%%%%%%%%%%%%%%%%


\begin{document}


\begin{center}
    % {\Huge \scshape Curriculum Vitae/Resume} \\ \vspace{1pt}
    {\Huge \scshape Gyanig Kumar} \\ \vspace{1pt}
    University of Colorado, Boulder \\ \vspace{1pt}
    \small \raisebox{-0.1\height}\faPhone\ +1-720-260-2451 ~ \href{mailto:gyanig.kumar@gmail.com}{\raisebox{-0.2\height
}\faEnvelope\ \underline{gyanig.kumar@gmail.com}} ~ 
    \href{https://linkedin.com/in//}{\raisebox{-0.2\height}\faLinkedin\ \underline{{\href{https://www.linkedin.com/in/gyanig-k-b36666146/}{https://www.linkedin.com/in/gyanig-k}}}}  ~
    \href{https://github.com/}{\raisebox{-0.2\height}\faGithub\ \underline{https://gyanigk.github.io}}
    \href{https://scholar.google.com/citations?user=g1kxkMUAAAAJ&hl=en}{\raisebox{-0.2\height}\aiGoogleScholar\ \underline{Google Scholar}}
    \vspace{-8pt}
\end{center}


%-----------EDUCATION-----------

\section{Education}
  \resumeSubHeadingListStart
    \resumeSubheading
      {University of Colorado, Boulder}{Aug 2024 -- Jul 2026}
      {M.S. in Computer Science}{CGPA: 3.91/4.0}
      \resumeItemListStart
        \resumeItem{\textbf{Relevant Coursework}: Advanced Robotics, Machine Learning, Advanced Computer Vision}
      \resumeItemListEnd
    \resumeSubheading
      {KIIT, India}{Jul 2019 -- Jul 2023}
      {B.Tech in Computer Science and Systems Engineering}{CGPA: 3.88/4.0 (9.05/10)}
      \resumeItemListStart
        \resumeItem{\textbf{Relevant Coursework}: Data Structures, Artificial Intelligence, Computer Networks, Compiler Design}
      \resumeItemListEnd
  \resumeSubHeadingListEnd
% \vspace{-14pt}
%-----------PROGRAMMING SKILLS-----------
\section{Technical Skills}
\begin{itemize}[leftmargin=0.15in, label={}]
  \item{
    \textbf{Programming}: Python, C/C++, C#, Java, Bash \\
    \textbf{Frameworks \& Platforms}: TensorFlow, PyTorch, TensorFlow.js, LangChain, ChromaDB, Socket.io \\
    \textbf{Libraries}: NumPy, Pandas, Keras, Librosa, Kivy, OpenCV, Matplotlib \\
    \textbf{Tools \& Technologies}: ROS, AWS, GCP, Unity 3D, Git, Linux, Fusion 360, MATLAB \\
    \textbf{Expertise}: Diffusion based I2I Translation, Multi-modal learning in HRI, Gaze Estimation, \& Object Pose Estimation
  }
\end{itemize}

\vspace{-8pt}
\section{Research Interests}
% \cusemph{HackerRank}: Gold Level C++ programmer }\\
{Generative-AI, Multi-modal learning, 3D Gaze Estimation, \& Object Registration }

\vspace{-4pt}

%-----------EXPERIENCE-----------
\section{Experience}
\vspace{2pt}
  \resumeSubHeadingListStart
  \resumeSubheading
      {Course Manager, CSCI 3302 - Intro to Robotics}{Jan'25 -- Present}
      {Univ. of Colorado, Boulder}{}
      \resumeItemListStart
      \resumeItem{\textbf{Coordinated course logistics and student support} for 120 undergrad and graduate students,\\ managing assignments, Canvas content, technical troubleshooting, and office hours assistance.}
\resumeItem{\textbf{Oversaw grading and assessments} ensuring timely feedback, utilizing MOSS for plagiarism checks,\\ and upholding fairness in evaluations.}

       \resumeItem{\textbf{Manage grading and review} for a class of 120 undergrad and graduate students, ensuring timely grading, utilizing \\MOSS for plagiarism detection, and maintaining fairness in assessments}
     \resumeItemListEnd
    \resumeSubheading
      {CAIRO HCI Lab, Univ. of Colorado, Boulder}{Sept'24 -- Present}
      {Researcher under Prof. Bradley Hayes}{}
      \resumeItemListStart
        \resumeItem{ Conducting research on enhancing \textbf{human intent recognition} using \textbf{Partially Observable Markov Decision \\Processes (POMDP)}-based algorithms.}
        \resumeItem{Integrating advanced \textbf{vision-based learning} techniques with the \textbf{Sawyer Robotic Manipulator}\\ to improve human-robot interaction and decision-making efficiency.}
      \resumeItemListEnd
    
    \resumeSubheading
      {I3D HCI Lab, Indian Institute of Science, Bangalore}{April'22 -- Aug'24}
      {Machine Learning Intern, Research Assistant under Prof. Pradipta Biswas(Phd. Cambridge)}{}\\
      \resumeItemListStart
      \resumeItem{\textbf{Developed and validated a gaze estimation mapping module}, achieving \textbf{95\% mAP} accuracy in predicting user\\ intent in complex robotic pick-and-place task.}
      \resumeItem{Partnered with \textbf{Collins Aerospace} to develop a mixed reality-based assembly system using \textbf{custom-trained object detection models} with \textbf{tailored eye gaze and hand tracking} for efficient, real-time assembly guidance.}
      \resumeItem{Conducted a comprehensive study comparing GAN models for \textbf{generating realistic synthetic datasets \\(Simulation-to- Real)} achieving FID scores of \textbf{0.001, 50.574, and 63.000} across three datasets.}
      \resumeItem{\textbf{Developed a robust motion tracking framework} with a markerless ByteTrack model, achieving \textbf{97.8\% mAP}\\ detection accuracy in dynamic environments, including moving camera setups such as headset-mounted systems.}
      \resumeItem{Optimized the \textbf{I2D-Net Eye Gaze Estimation model}, reducing parameters by \textbf{0.5x} and \textbf{2x} inference speed,\\ enabling a relatively faster gaze tracking system for interactive automotive HUDs}
    \resumeItem{\textbf{Collaborated with Faurecia} to extend interactive automotive heads-up display features with a \\\textbf{gaze-based vs gesture-based} on-road distraction detection system.}
     \resumeItem{Published \textbf{5 research papers} at prestigious top conferences like \textbf{ICRA, IUI} and Journals like \textbf{JMUI}.}
      \resumeItemListEnd
    
    \resumeSubheading
      {Part-Time CS Teacher, Polytechnic College Suriname, Suriname, SA}{March'23 -- Sept'24}
      {Data Science Program}{\href{https://drive.google.com/file/d/1M6aUJF2gDGg3sifKaebum_3h_qnQclOr/view?usp=sharing}{\underline{Appointment Letter}}}
      \resumeItemListStart
        \resumeItem{Conducted classes, designed test \& assignments, and evaluated 50 enrolled students for \textbf{University-level}\\\textbf{Mathematics II and Statistics I Courses} for the semester with 15 hours avg. per week.}
        \resumeItem{Undertook \textbf{50+ students} for guidance in Programming \& Prototyping in a hackathon.}
        
    \resumeItemListEnd
    
  \resumeSubHeadingListEnd
\vspace{-12pt}


\section{Projects}
  \resumeSubHeadingListStart
      \resumeSubheading
    {Ansh.AI: Conversational Diary Companion}{Mar'25}
    {HackCU11 $|$ \href{https://github.com/sakshamk6999/HackCU11}{GitHub} $|$ \href{https://youtu.be/1wwqeeN3aRM}{Demo}}{}
    \resumeItemListStart
      \resumeItem{Built an empathetic AI for self-reflection using \textbf{LangChain and HeyGen} for natural conversations.}
      \resumeItem{Enabled real-time memory with \textbf{ChromaDB for personalized insights} and pattern recognition.}
      \resumeItem{Developed with Python, Flask, React, and \textbf{AMD Minisforum AI 370 for NPU+GPU compute}.}
      \resumeItem{Secured \textbf{2nd Position} in the AI Track of the Hackathon}
    \resumeItemListEnd
    \resumeSubheading
      {Autonomous Driving Challenge with AWS DeepRacer}{Aug'24 -- Dec'24}
      {University of Colorado, Boulder $|$ \href{https://github.com/gyanigk/Advanced-Robotics/tree/Final-Project/FINAL_PROJECT}{\underline{GitHub}}}{}
      \resumeItemListStart
        \resumeItem{Built a 1/18-scale autonomous vehicle using \textbf{ROS 2 Humble} on Raspberry Pi 4, navigating via \textbf{RGB camera-based} blue line following and \textbf{LiDAR odometry}.}
        \resumeItem{Applied \textbf{HSV thresholding} and \textbf{contour detection} with OpenCV for traffic sign detection, achieving \textbf{~80\% stop sign accuracy}, and implemented \textbf{EKF SLAM} for robust localization.}
        \resumeItem{Designed a modular ROS 2 system with nodes for camera processing, line detection, PID control, and \textbf{slam\_toolbox} for real-time mapping and \textbf{Tkinter}-based telemetry visualisation.}
    \resumeItemListEnd
    \resumeSubheading
      {Audio-Visual Correspondance Task using Few Shot Learning}{Dec'21 - June'22}
      {Supervisor: Dr. Rajdeep Chatterjee}{}\\
      \resumeItemListStart
          \resumeItem{Successfully reproduced the results from the \href{https://arxiv.org/abs/1705.08168}{\textit{\underline{"Look Listen \& Learn"}}} paper, implementing the L3Net model.}
        \resumeItem{Curated a specialized subset of \textbf{34 classes from the KINETIC600 dataset} to train a modified Zero-Shot Learning architecture with an RCNN backbone, incorporating dual subnetworks: \textbf{Audio Encoder} and \textbf{Vision Encoder}.}
        \resumeItem{Employed advanced feature extraction techniques such as \textit{FFT} and \textit{Mel-Filter Bank spectrum} for audio features, aligning them with corresponding image samples for improved model performance.}
        \resumeItem{Achieved \textbf{0.180 (Vision)} and \textbf{0.245 (Audio)} accuracy, demonstrating significant potential for further optimization and performance enhancement.}
        \resumeItemListEnd
      \vspace{-0.5pt}
  \resumeSubHeadingListEnd
\vspace{-14pt}

\section{Other Projects}
  \resumeSubHeadingListStart
    \resumeSubheading
      {Lox: Interpreter for Python-like language $|$ Compiler Design, OOPs $|$ \href{https://github.com/gyanigk/Lox}{\underline{Github}}}{Sept'20-Dec'20}
      {}{}
      \vspace{-18pt}
      \resumeItemListStart
        \resumeItem{Built an \textbf{end-to-end interpreter} that takes Python-like syntax based on \href{https://craftinginterpreters.com/}{\underline{Crafting Interpretors}}.}
        \resumeItem{Implemented features such as \textit{parsing}, \textit{control flow}, \textit{hashes}, \textit{garbage collection}, \textit{superclasses} etc.}
      \resumeItemListEnd

    \resumeSubheading
      {V-BIKES: Deploying Enhanced Safety Measures in two-wheeler vehicles $|$ \href{https://drive.google.com/file/d/1kB5Lr3pf-EXvssMpe6t8tkbkbpHlXwFj/view?usp=share_link}{\underline{Award}}}{Sept'17-Feb'18}
      {}{}
      \vspace{-18pt}
      \resumeItemListStart
        \resumeItem{Based on the field of robotics, we developed an \textbf{array of non-evasive sensors for collision detection} in a helmet}
        \resumeItem{Built End-to-End support for Immediate \textbf{Ambulance Alert \& injury Profile} using ESP32 Kit.}
        \resumeItem{{\href{https://drive.google.com/file/d/1JOfyKVWSockjHTPi5B7Y7iNZmLMB-D6l/view?usp=share_link}{Awarded as \textbf{\underline{best in Smart Mobility Project in Atal Marathon}}  by Government of India \\among 1200+ projects and \textbf{3rd Position} in \textbf{C.S.I.R. Innovation Awards.}}} \\}
        
        
      \resumeItemListEnd
    
  \resumeSubHeadingListEnd
\vspace{-14pt}

\section{Publications}
\begin{enumerate}
        
    \item {L.R.D. Murthy, \textbf{Gyanig Kumar}, Pradipta Biswas, M. Madan, S. Deshmukh\\
    \textit{"Efficient Interaction with Automotive Heads-Up Displays using Appearance-based Gaze Tracking"}\\
    \textbf{Published} in \textit{Work-in-Progress Track of the 14th International Conference on Automotive User Interfaces and Interactive Vehicular Applications (AutomotiveUI 2022)}\\
    \href{https://dl.acm.org/doi/pdf/10.1145/3544999.3554818}{\underline{ACM Digital Library Link}} | \href{https://www.youtube.com/watch?v=lqWYUUv88_U}{\underline{Video Submission}}}
    
    \item {Yash Kumar Sahu, A. Mukhopadhyay, \textbf{Gyanig Kumar}, Ashok Kumar, Pradipta Biswas\\
    \textit{"A Comparative Study on Image Translation GAN Models to Improve Object Detection in Low-Resource Domains"}\\
    \textbf{Published} in \textit{2024 International Conference on Vehicular Technology and Transportation Systems (ICVTTS)}\\
    \href{https://ieeexplore.ieee.org/document/10763945}{\underline{IEEE Xplore Link}}}
    
    \item {Mukund Mitra, A.A. Patil, G. Mothish, \textbf{Gyanig Kumar}, A. Mukhopadhyay, L.R.D. Murthy, P.P. Chakraborty, Pradipta Biswas\\
    \textit{"Multimodal Target Prediction for Rapid Human-Robot Interaction"}\\
    \textbf{Published} in \textit{29th ACM Conference on Intelligent User Interfaces (ACM IUI) 2024}\\
    \href{https://dl.acm.org/doi/abs/10.1145/3640544.3645229}{\underline{ACM Digital Library Link}}}
    
    \item {Mukund Mitra, \textbf{Gyanig Kumar}, P.P. Chakraborty, Pradipta Biswas\\
    \textit{"Enhanced Human-Robot Collaboration with Intent Prediction using Deep-IRL"}\\
    \textbf{Published} in \textit{IEEE International Conference on Robotics and Automation (ICRA) 2024}\\
    \href{https://ieeexplore.ieee.org/abstract/document/10610595}{\underline{IEEE Xplore Link}}}
    
    \item {Subin Raj, L.R.D. Murthy, T.A. Shanmugam, \textbf{Gyanig Kumar}, A. Chakrabarti, Pradipta Biswas\\
    \textit{"Mixed Reality and Deep Learning-based System for Assisting Assembly Processes"}\\
    \textbf{Published} in \textit{Journal on Multimodal User Interfaces (2023)}\\
    \href{https://link.springer.com/article/10.1007/s12193-023-00428-3}{\underline{Springer Link}}}


\end{enumerate}

%-----------PROJECTS-----------
\section{Academic Projects}
    \vspace{-5pt}
    \resumeSubHeadingListStart
      \resumeProjectHeading
          {\textbf{Chat rooms with multilingual conversation support } $|$ \emph{Supervisor: Prof. Bindu Agarwalla}}{Feb'22 - April'22}
          \resumeItemListStart
            \resumeItem{Novel introduction of \textbf{translating any incoming messages} on a chat platform.}
            \resumeItem{Intricate Web server application using Express.js \& Socket.io for creating \textbf{bi-directional}\\ \textbf{messaging }passing with efficient query handling subroutine.}
            \resumeItem{\textbf{Google Translate API} provides faster translation and dynamic language support}
          \resumeItemListEnd
          \vspace{-13pt}
\vspace{-8pt}



\section{Volunteering}
  \resumeSubHeadingListStart
    \resumeSubheading
      {Teaching Assistant, Konnexions Society, KIIT}{Nov'21 - April'22}
      {}{}
      \vspace{-18pt}
      \resumeItemListStart
        \resumeItem{Teaching Assistant to a class over 100+ Students. | \href{https://drive.google.com/file/d/1FdKyJ08dQoBWGHq5wX4FWgagSQDSOkkY/view?usp=sharing}{\underline{Appreciation Letter}}}
        \resumeItem{Managed a No-Credit Course focused on \textbf{Mathematics in Machine Learning} for a semester.}
      \resumeItemListEnd

    \resumeSubheading
    {Student Member, Google Developer Student Community, KIIT}{Oct'21 - Oct'22}
      {}{}
      \vspace{-18pt}
      \resumeItemListStart
        \resumeItem{Contributed in Organizing Online Seminar, Hackathons \& Workshops.}
        \resumeItem{Worked with the ML team on a \textit{``PetFinder.my - Pawpularity Research Contest"} on Kaggle, using ensemble with \textbf{custom Vision models.} (e.g. BiT ,ViT, EffNetB2-B5 )}
        \resumeItem{Secured \textbf{2\textsuperscript{nd} Position} out of 50 teams in KIIT university.}
        \resumeItemListEnd
    \resumeSubheading
    {Lead Data Acquisition Engineer, Jaggurnaut Racing, KIIT}{Feb'22 - Feb'23}
      {}{}
      \vspace{-18pt}
      \resumeItemListStart
        \resumeItem{Participated in BAJA SAE India 2022 ATV design and Building Competition}
        \resumeItem{Build a dashboard system for ATV with sensor data information}
        \resumeItem{Secured \textbf{1\textsuperscript{st} Position} in KIIT University and \textbf{16\textsuperscript{th}} among 100+ teams in India}
    \resumeItemListEnd
  \resumeSubHeadingListEnd
\vspace{-14pt}

% % -----------Certificates-----------
\section{Certifications}
 \begin{itemize}[leftmargin=0.15in, label={}]
    {\item{
     {\href{https://coursera.org/share/6ee40d5a2d01fcabbafe4e0fdf5a743e}{\textbf{\underline{TensorFlow DeepLearning Specilization}}} ,
     \textit{DeepLearning.ai} - Jan'22\\
     
     \href{https://drive.google.com/file/d/15z48oaRECDBw-B6Nh2YyCx1UZtYELvQo/view?usp=share_link}{\textbf{\underline{5th Summer School on AI}}} ,
     \textit{ CVIT, IIIT Hyderabad} - Sept'21 \\

     \href{https://coursera.org/share/6178c57f2559a529439c376c405bf640}{\textbf{\underline{Coursera course on Machine Learning}}} ,
     \textit{ Stanford} - July'20 \\
     
    }}
 \end{itemize}
 
%-----------Achievements-----------
\section{\href{https://drive.google.com/drive/folders/1TIIW4SO32AS2VVs_9XBNl7Uw9k_bjr-y?usp=share_link}{Achievements}}
 \begin{itemize}[leftmargin=0.15in, label={}]
    % {\item{
    %  {\href{https://www.twitch.tv/p4r7_}{\textbf{\underline{Official Twitch Affliated Streammer}}} under Science \& Technology. - 2020} \\
    %  {Participated in Water Rover Competition \textbf{IIT BOMBAY Techfest}. - 2019\\
    % }
    \item{
        Secured \textbf{2nd Place} in HACKCU11, AI Track. - [2025] \\
        Secured \textbf{2nd Place} in AWS JAM Hackathon. - [2025]
    }
 \end{itemize}


\section{Hobbies}
\cusemph{HackerRank}: Gold Level C++ programmer }\\
{Soccer, Programming, Electric Guitar, Composing Music}

\end{document}
